\section*{Chapter \ref{ch:b2:trees-and-clustering} --- Learning from Examples: Decision Trees and Clustering}

\begin{solution}{exo:b2:trees-and-clustering:1}
\omterm{def:b2:trees-and-clustering:confusion}{Accuracy}: $(224 + 35)/300 = 0.863$. Late loans caught: $35/67 \approx 0.522$. Right ``late'' predictions: $35/(35 + 9)
\approx 0.795$. If only false alarms cost anything, the library should watch the last number, the fraction of alarms that
are right, and could raise the bar for calling a loan late; \omterm{def:b2:trees-and-clustering:confusion}{accuracy} mixes the two kinds of error with equal weight.
\end{solution}

\begin{solution}{exo:b2:trees-and-clustering:2}
\omoutput{code/bachelor-2/29-trees-and-clustering/outputs/exo_scaling.txt}{Raw \omterm{def:b2:trees-and-clustering:examples}{features} against scaled ones.}
The distance in metres ranges over almost 6,000 and the pages over about 800, while the late returns range over 6: in the
raw distance, a difference of one late return weighs as much as one metre. The neighbours are then the members who live
at about the same distance and borrow books of about the same length, whatever their late returns, and one of the two
\omterm{def:b2:trees-and-clustering:examples}{features} of the rule is lost. Scaling every \omterm{def:b2:trees-and-clustering:examples}{feature} to $[0, 1]$ gives each the same weight.
\end{solution}

\begin{solution}{exo:b2:trees-and-clustering:3}
The function visits a node, reports its point if it lies in the rectangle, and goes down to the left only if the
rectangle reaches coordinates at most the node's on its axis ($\mathit{lo}_a \leq p_a$), to the right only if it reaches
coordinates at least the node's ($\mathit{hi}_a \geq p_a$).

\omoutput{code/bachelor-2/29-trees-and-clustering/outputs/exo_range.txt}{Three rectangles in a tree of 10,000 points.}
The nodes visited are the points reported plus the nodes whose region only crosses the rectangle's border: about twice
the answer for the square, and 13 for an empty point query, a path down the tree. A long thin strip crosses many
regions, so it visits more nodes per point reported.
\end{solution}

\begin{solution}{exo:b2:trees-and-clustering:4}
The whole set has $9/24 = 0.375$ late, so $H = h(0.375) \approx 0.954$ bits. The left part has $6/8 = 0.75$ late,
$h(0.75) \approx 0.811$; the right part $3/16 = 0.1875$, $h(0.1875) \approx 0.696$. The parts weigh $8/24 = \tfrac13$ and
$16/24 = \tfrac23$, so the gain is $0.954 - \tfrac{1}{3} \cdot 0.811 - \tfrac{2}{3} \cdot 0.696 \approx 0.220$ bits. On
the \omterm{def:b2:trees-and-clustering:entropy}{entropy} figure, the average of the parts is the point of the chord between them above $p = 0.375$, two thirds of
the way from $0.75$ to $0.1875$, not its middle.
\end{solution}

\begin{solution}{exo:b2:trees-and-clustering:5}
\omoutput{code/bachelor-2/29-trees-and-clustering/outputs/exo_lloyd_step.txt}{One iteration, then the assignment that
stops the algorithm.}
From A, E and I, the points A, B, C go to A; D, E, F, G, H to E (D is $2.26$ from E and $3.65$ from A); I, J, K, L to I.
The cost is $42.54$. The \omterm{def:b2:trees-and-clustering:clustering}{centroids} are $(1.50, 1.77)$, $(7.28, 1.94)$ and $(4.65, 8.00)$, and the same assignment now
costs $28.78$. With these centres no point changes cluster, so the algorithm stops after one iteration, with D in the
right-hand group.
\end{solution}

\begin{solution}{exo:b2:trees-and-clustering:6}
\omoutput{code/bachelor-2/29-trees-and-clustering/outputs/exo_identifier.txt}{One branch per value: the loan's number wins.}
A split with one branch per loan leaves a single loan in each branch, so every branch has \omterm{def:b2:trees-and-clustering:entropy}{entropy} $0$ and the gain is the
whole \omterm{def:b2:trees-and-clustering:entropy}{entropy} of the labels, the largest possible. The tree would ask ``which loan is it?'' and learn nothing that applies
to a new loan, whose number is new. The gain ratio divides by the \omterm{def:b2:trees-and-clustering:entropy}{entropy} of the branch sizes, $\log_2 700 \approx
9.45$ bits for the number, and lowers its score to $0.093$, but that is still above the $0.064$ of the late returns: here
it is not enough. Identifiers must be left out of the \omterm{def:b2:trees-and-clustering:examples}{features}; binary splits limit the damage, since one threshold on the
number gains almost nothing.
\end{solution}

\begin{solution}{exo:b2:trees-and-clustering:7}
\omoutput{code/bachelor-2/29-trees-and-clustering/outputs/exo_crossval.txt}{The mean of the five folds for each depth.}
Cross-validation chooses depth 2, with a mean of $0.846$, as the \omterm{def:b2:trees-and-clustering:examples}{test set} did in \cref{out:b2:trees-and-clustering:overfit},
without looking at it. The five folds differ by up to $0.11$ for the same depth ($0.07$ at depth 2): one validation set of 140 is a noisy
judge, and the mean of five is steadier.
\end{solution}

\begin{solution}{exo:b2:trees-and-clustering:8}
Write $a_t$ for the $t$-th assignment, $c_t$ for the centres from which it is computed, and $\mathrm{cost}(c, a)$ for
the $k$-means cost. The assignment step gives each point its nearest centre, so $\mathrm{cost}(c_t, a_t) \leq
\mathrm{cost}(c_t, a_{t-1})$. The \omterm{def:b2:trees-and-clustering:clustering}{centroid} step replaces the centre of each non-empty cluster of $a_{t-1}$ by its
\omterm{def:b2:trees-and-clustering:clustering}{centroid}, which by \cref{lem:b2:trees-and-clustering:centroid} does not increase that cluster's sum, and leaves empty
clusters' centres in place; so $\mathrm{cost}(c_t, a_{t-1}) \leq \mathrm{cost}(c_{t-1}, a_{t-1})$. Together,
$\mathrm{cost}(c_t, a_t) \leq \mathrm{cost}(c_{t-1}, a_{t-1})$: the recorded cost never increases.

For termination, let $D(a)$ be the cost of an assignment $a$ with each centre at the \omterm{def:b2:trees-and-clustering:clustering}{centroid} of its cluster, empty
clusters contributing nothing; then $D(a_{t-1}) = \mathrm{cost}(c_t, a_{t-1})$, since the \omterm{def:b2:trees-and-clustering:clustering}{centroid} step computes $c_t$
from $a_{t-1}$. The two steps and the lemma give $D(a_{t-1}) = \mathrm{cost}(c_t, a_{t-1}) \geq \mathrm{cost}(c_t, a_t)
\geq D(a_t)$. If $D(a_t) = D(a_{t-1})$, the last inequality is an equality, and by the uniqueness in the lemma each
centre of $c_t$ with a non-empty cluster in $a_t$ is already its \omterm{def:b2:trees-and-clustering:clustering}{centroid}; the others stay in place. So the next centres
are $c_t$ again, the next assignment is $a_t$ again, and the algorithm stops. Hence, as long as it runs, $D$ strictly
decreases, so no assignment comes back; there are at most $k^n$ assignments, so the algorithm stops.
\end{solution}

\begin{solution}{exo:b2:trees-and-clustering:9}
\omoutput{code/bachelor-2/29-trees-and-clustering/outputs/exo_complete.txt}{The eleven merges of complete \omterm{def:b2:trees-and-clustering:hac}{linkage}.}
The first five merges are A with B ($1.26$), F with G ($1.58$), J with K ($1.77$), C with AB ($1.92$: its farthest point
in AB is A) and E with FG ($2.21$). Complete \omterm{def:b2:trees-and-clustering:hac}{linkage} measures a cluster by its farthest point: D is $3.65$ from A, the
farthest of A, B, C, but $4.46$ from G, the farthest of E, F, G, H, so it joins A, B, C. Single \omterm{def:b2:trees-and-clustering:hac}{linkage} measures by the
nearest point: D is $2.26$ from E but $2.41$ from B, so it joins E's group.
\end{solution}

\begin{solution}{exo:b2:trees-and-clustering:10}
Run Kruskal's algorithm on the complete graph: it examines the edges by increasing length and accepts an edge when its
ends lie in different components. We show by induction on $m$ that after $m$ merges, the clusters of single \omterm{def:b2:trees-and-clustering:hac}{linkage} are
the components after $m$ accepted edges. At $m = 0$ both are the singletons.

Suppose it holds after $m$ merges. An edge whose ends are now in different components was never examined: when it was,
its ends were in different components, since components only grow, and it would have been accepted, putting them in the
same component. So the next edge Kruskal accepts is the shortest edge between two different components. The next merge of
single \omterm{def:b2:trees-and-clustering:hac}{linkage} joins the two clusters at the least single \omterm{def:b2:trees-and-clustering:hac}{linkage}, the least distance between a point of one and a
point of the other: it is the shortest edge between two different clusters. Since the distances are distinct, it is
the same edge, and both join the same two sets.

After $n - k$ merges there are $k$ clusters, the components after $n - k$ accepted edges. With distinct weights the
minimum spanning tree is unique and is Kruskal's, made of its $n - 1$ accepted edges in increasing order; the
components after the first $n - k$ are those of the tree without the last $k - 1$, its $k - 1$ heaviest edges.
\end{solution}

\begin{solution}{exo:b2:trees-and-clustering:11}
\omoutput{code/bachelor-2/29-trees-and-clustering/outputs/exo_kmeanspp.txt}{100 runs from each seeding, costs rounded.}
Uniform seeding finds the best \omterm{def:b2:trees-and-clustering:clustering}{clustering}, cost 131, in 53 runs of 100, and stops at 17 different costs; $k$-means++
finds it in 80 runs and stops at 12 different costs. With five groups, a uniform start leaves some group without a centre
with fair probability; the $k$-means++ start, by favouring points far from the centres chosen, rarely does. Neither
avoids every local minimum, so keeping the best of several runs remains the practice.
\end{solution}

\begin{solution}{exo:b2:trees-and-clustering:12}
Put the points on a circle centred on the query, for example the 20 integer points with $x^2 + y^2 = 625$, and the
query at the origin. Every point is at squared distance 625, so the best distance never improves after the first point,
and a subtree is skipped only when $(q_a - p_a)^2 \geq 625$. Every coordinate of every point lies in $[-25, 25]$, so the
test fails except at the points with a coordinate $\pm 25$ on the node's axis, whose far side holds nothing on the
circle.

\omoutput{code/bachelor-2/29-trees-and-clustering/outputs/exo_kd_worst.txt}{The query at the centre visits all 20 nodes;
moving it a little barely helps.}
The search visits all 20 nodes, the cost of the brute force plus the overhead of the tree. The pruning needs the best
distance to be small compared with the distances to the splitting lines; points all at nearly the same distance from
the query defeat it.
\end{solution}

\begin{solution}{pb:b2:trees-and-clustering:1}
\textbf{1.} The loans might come in an order that matters, by date for example; without shuffling, the \omterm{def:b2:trees-and-clustering:examples}{test set} could
come from another period than the \omterm{def:b2:trees-and-clustering:examples}{training set}, and measure something else.

\textbf{2.} The \omterm{def:b2:trees-and-clustering:examples}{training set} grows the trees, the validation set chooses the depth, the \omterm{def:b2:trees-and-clustering:examples}{test set} measures the chosen
tree. A \omterm{def:b2:trees-and-clustering:examples}{test set} used to choose anything becomes part of the training, and its score is then optimistic.

\textbf{3.} The two fractions come from the same source and differ by chance: on 300 loans the standard deviation of a
fraction near the expected $0.27$ is about $0.026$, and the \omterm{def:b2:trees-and-clustering:examples}{test set}'s $0.217$ is two of them below it. A stratified split, which keeps the same
proportion of late loans in each part, removes that source of variation.

\textbf{4.} Even the planted rule is wrong on the flipped labels, one in ten on average, so about $0.9$, give or take
the chance of a single set.

\textbf{5.} Always answering ``on time'' scores $1 - 65/300 \approx 0.783$. A \omterm{def:b2:trees-and-clustering:learning}{classifier} that does not beat it has
learned nothing useful.

\textbf{6.} One would guess the late returns, and perhaps the distance or the days lent; nothing in the data says so in
advance.

\textbf{7.} Depth 2, with $0.885$ on validation; depths 3 and 4 score $0.880$, slightly less. Among nearly equal scores
the shallowest tree is the simplest, and the least likely to have fitted noise.

\textbf{8.} A loan whose member has more than 1.5 late returns, that is at least 2, and lives more than 2,495 m away.

\textbf{9.} The planted rule is ``at least 2 late returns and at least 2,500 m''. The distances are multiples of 10, so
``more than 2,495'' and ``at least 2,500'' agree on every possible loan, and the late returns are whole numbers: the
tree computes exactly the rule, and agrees with it on every test loan.

\textbf{10.} On the noisy training loans, the left branch has 69 late labels out of 457, all of them noise, and they are
not spread evenly: 34 of 176 among the short books, 35 of 281 among the long ones. Splitting there lowers the \omterm{def:b2:trees-and-clustering:entropy}{entropy} a
little, so the gain is positive, but both parts keep a majority on time and the predictions do not change. Pruning,
which removes a split whose two leaves agree, would remove it.

\textbf{11.} The 700 loans have $34 + 35 + 5 + 138 = 212$ late, $H \approx 0.885$ bits. The left part has $69$ late of
$457$, \omterm{def:b2:trees-and-clustering:entropy}{entropy} about $0.612$; the right part $143$ of $243$, about $0.977$. The gain is $0.885 - \tfrac{457}{700} \cdot
0.612 - \tfrac{243}{700} \cdot 0.977 \approx 0.146$ bits.

\textbf{12.} The tree recovers the planted rule at depth 2: test \omterm{def:b2:trees-and-clustering:confusion}{accuracy} $0.923$, against a majority baseline of
$0.783$.

\textbf{13.} Below depth 2, the tree has not learned the rule; beyond, the only patterns left in the training labels are
the flips, which are random, so each new split fits flips that new loans do not share.

\textbf{14.} Once the depth is chosen, the validation loans are only more data: the tree grown on 700 loans estimates
the thresholds and the majorities better.

\textbf{15.} Yes. On 300 loans with labels wrong with probability $0.1$, the number of wrong labels varies: its standard
deviation is $\sqrt{300 \cdot 0.1 \cdot 0.9} \approx 5.2$, so this \omterm{def:b2:trees-and-clustering:examples}{test set} had about 23 flips instead of 30. A single
\omterm{def:b2:trees-and-clustering:examples}{test set} of 300 gives \omterm{def:b2:trees-and-clustering:confusion}{accuracy} to within a few hundredths only.

\textbf{16.} On the loans with 15~\% noise, $k$-nearest neighbours reaches about $0.84$ with all five \omterm{def:b2:trees-and-clustering:examples}{features} against
$0.863$ for the depth-2 tree, because the three useless \omterm{def:b2:trees-and-clustering:examples}{features} enter every distance; with the two right \omterm{def:b2:trees-and-clustering:examples}{features} it
does about as well as the tree. It also keeps all the \omterm{def:b2:trees-and-clustering:examples}{examples}, where the tree keeps two questions.

\textbf{17.} The flips would hide the rule more: validation would choose a shallow tree, still the rule if the data are
enough, with an \omterm{def:b2:trees-and-clustering:confusion}{accuracy} of at most about $0.7$; with fewer loans it might choose depth 0 or 1.

\textbf{18.} Rerun the whole procedure on other shuffles of the data, or on samples drawn with replacement, and compare
the \omterm{def:b2:trees-and-clustering:examples}{features} and thresholds chosen: a rule that keeps the same two questions with nearby thresholds is stable.

\textbf{19.} The members have no label: nobody has said which group each belongs to. The algorithm has only the points.

\textbf{20.} $k = 3$: the cost falls by a factor of four from 1 to 2 and more than three from 2 to 3, then by less than a
quarter at each step. Past the bend, the extra clusters only split real groups.

\textbf{21.} About 31 loans a year kept 9 days: frequent readers of short books. About 12 a year kept 24 days: readers who
keep their books long. About 4 or 5 a year kept 13 days: occasional borrowers.

\textbf{22.} \omterm{met:b2:trees-and-clustering:lloyd}{Lloyd's algorithm} stops at a local minimum that depends on the start; the best of 10 starts is much more
likely to be the best \omterm{def:b2:trees-and-clustering:clustering}{clustering}, as in \cref{out:b2:trees-and-clustering:kmeanspp}.

\textbf{23.} Both \omterm{def:b2:trees-and-clustering:examples}{features} range over a few tens, so scaling would change little; in general, as with $k$-nearest
neighbours, a \omterm{def:b2:trees-and-clustering:examples}{feature} with a larger range would decide the distances alone.

\textbf{24.} \omterm{def:b2:trees-and-clustering:hac}{Hierarchical agglomerative clustering}: its \omterm{def:b2:trees-and-clustering:hac}{dendrogram} shows every number of clusters at once, and a long
vertical gap between two merge heights suggests where to cut.
\end{solution}
